\section*{Chapitre \ref{ch:ml:clustering-regimes-and-anomaly-detection} --- Partitionnement (clustering), régimes et détection d'anomalies}

\begin{solution}{exo:ml:clustering-regimes-and-anomaly-detection:1}
$\sqrt{2(1 - 0.5)} = 1$. Des titres parfaitement corrélés sont à distance 0, des titres non corrélés à $\sqrt2 = 1.41$ (et parfaitement
anticorrélés à 2).
\end{solution}

\begin{solution}{exo:ml:clustering-regimes-and-anomaly-detection:2}
Durées $1/(1 - p)$~: 82 jours calmes et 24 jours de tension (la matrice ajustée, avec 0.988 et 0.958 arrondis, donne 82.2 et 23.9). La
part de long terme des jours de tension est $P_{12}/(P_{12} + P_{21}) = 0.225$, contre 0.215 sur les années de test.
\end{solution}

\begin{solution}{exo:ml:clustering-regimes-and-anomaly-detection:3}
115 jours-compte légitimes sont signalés à 1~\%~; avec les 100 cas réels également signalés, la précision vaut $100/215 = 0.47$.
\end{solution}

\begin{solution}{exo:ml:clustering-regimes-and-anomaly-detection:4}
Le filtre n'utilise que les données jusqu'à ce jour, et la veille d'un basculement est un jour calme~: son rendement est tiré de la
distribution. Un basculement est annoncé par les jours de tension qui le suivent, que le filtre n'a pas encore vus~; sa probabilité
ce jour-là est simplement celle d'un basculement un jour calme, d'environ 1~\%, infléchie par le rendement du jour.
\end{solution}

\begin{solution}{exo:ml:clustering-regimes-and-anomaly-detection:5}
Chaque action est exposée au marché, si bien que les corrélations brutes sont dominées par un facteur commun~: tous les titres se ressemblent, et le
\omterm{def:ml:clustering-regimes-and-anomaly-detection:cluster}{partitionnement} retient des différences de bêta ou de volatilité spécifique plutôt que des secteurs. Retirer le marché laisse les
corrélations sectorielles déterminer les distances.
\end{solution}

\begin{solution}{exo:ml:clustering-regimes-and-anomaly-detection:6}
Les probabilités sont lissées, ou le modèle a été estimé sur tout l'échantillon, drawdowns compris~: ses états et ses paramètres savent
où se trouvaient les drawdowns, et la probabilité lissée augmente avant chaque basculement par construction (0.35 contre 0.02 les jours calmes
ici). Il faut relancer le modèle comme filtre, réestimé uniquement sur les données antérieures à chaque date, puis examiner à nouveau le résultat.
\end{solution}

\begin{solution}{exo:ml:clustering-regimes-and-anomaly-detection:7}
La log-vraisemblance sur les années de test vaut 7\,866.6 sous le modèle à deux états et 7\,864.4 sous celui à trois états~: l'état supplémentaire,
qui améliorait de 5.7 l'ajustement dans l'échantillon, en perd 2.2 hors échantillon. Le verdict du critère d'information tient.
\end{solution}

\begin{solution}{exo:ml:clustering-regimes-and-anomaly-detection:8}
$\alpha_t(j) = p(x_1,\dots,x_t, s_t = j) = \sum_ip(x_1,\dots,x_{t-1}, s_{t-1} = i)\P(s_t = j\mid s_{t-1} = i)p(x_t\mid s_t = j)$,
en utilisant la propriété de Markov de l'état et l'indépendance conditionnelle de $x_t$ sachant $s_t$~; soit
$\big(\sum_i\alpha_{t-1}(i)P_{ij}\big)b_j(x_t)$. En divisant par $\sum_j\alpha_t(j) = p(x_1,\dots,x_t)$, on obtient $\P(s_t = j\mid
x_1,\dots,x_t)$, la probabilité filtrée~; le code normalise à chaque étape, ce qui évite aussi les sous-dépassements numériques.
\end{solution}

\begin{solution}{pb:ml:clustering-regimes-and-anomaly-detection:1}
\textbf{Partie I.}
\begin{enumerate}
\item Avec les secteurs à 15~\%~: indice de Rand ajusté de 0.38 à 0.42 sur trois mois, de 0.60 à 0.75 sur un an, de 0.99 à 1.00 sur cinq
ans. À 8~\%~: de 0.03 à 0.06, 0.06, et de 0.15 à 0.21.
\item Accord sur demi-échantillon sur un an~: de 0.32 à 0.35 (secteurs marqués), de 0.06 à 0.08 (faibles). Sur des données réelles, il n'y a pas
d'étiquettes insérées~; la stabilité est le seul contrôle.
\item Regroupements pour les rapports de risque et parité de risque hiérarchique (HRP) lorsqu'ils sont stables~; non comme signal de trading, et non lorsqu'ils ne
survivent pas à un découpage.
\item Le \omterm{def:ml:clustering-regimes-and-anomaly-detection:cluster}{partitionnement hiérarchique} sur la distance de corrélation est ici légèrement meilleur et donne un arbre qui montre comment les groupes
s'emboîtent~; les \omterm{def:ml:clustering-regimes-and-anomaly-detection:cluster}{k-moyennes} (k-means) exigent un nombre fixe de groupes et une représentation euclidienne.
\end{enumerate}
\textbf{Partie II.}
\begin{enumerate}[resume]
\item Volatilité en régime calme de 0.80~\% et en régime de tension de 2.60~\% (valeurs insérées 0.8~\% et 2.5~\%), probabilités de maintien de 0.988 et 0.958 (0.99 et
0.95).
\item Le troisième état ajoute 5.7 à la log-vraisemblance dans l'échantillon pour six paramètres, en deçà de la valeur 23.5 du critère d'information, et
perd hors échantillon (exercice~7).
\item Filtrage 95.9~\%, lissage 97.7~\%, Viterbi 97.5~\%, seuil de volatilité 89.8~\%.
\item Médiane d'un jour pour le filtre (moyenne 1.5), nulle pour le lisseur (moyenne 0.8)~; la moyenne du seuil de volatilité est de 3.8 et il
manque cinq des 20 épisodes.
\end{enumerate}
\textbf{Partie III.}
\begin{enumerate}[resume]
\item Filtrage 0.032 (moyenne des jours calmes 0.041)~; lissage 0.353 (moyenne des jours calmes 0.019).
\item Ratio de Sharpe de 1.21 avec la probabilité lissée contre 1.00 avec la probabilité filtrée (achat et conservation 0.91).
\item À partir des probabilités filtrées, avec des paramètres estimés sur les données antérieures à chaque date.
\item Dans toute statistique calculée avec un ajustement sur tout l'échantillon~: normalisations, étiquettes de régime, retrait de tendance, nettoyage des valeurs aberrantes, et
lisseurs de Kalman utilisés là où un filtre était visé (livre~4, chapitre~19).
\end{enumerate}
\textbf{Partie IV.}
\begin{enumerate}[resume]
\item \emph{Dans quel régime sommes-nous~?} Le régime filtré a raison sur 95.9~\% des jours et repère un basculement vers la tension avec un délai médian
d'un jour~; le régime lissé a raison sur 97.7~\% des jours et le repère le jour même, parce qu'il connaît le futur. Au taux de 1~\% de faux
positifs, la précision de la \omterm{def:ml:clustering-regimes-and-anomaly-detection:anomaly}{forêt d'isolement} (isolation forest) est de 0.47 (rappel 1.00), celle de la meilleure règle de 0.40 (0.77) et celle de l'\omterm{def:ml:cross-sectional-deep-models:ae}{autoencodeur} de 0.22 (0.32).
\item Décrire l'état présent pour le dimensionnement du risque et pour conditionner d'autres modèles, vite et honnêtement, sans prédire
les basculements.
\item Non~: en second filet, avec revue, les règles étant les détecteurs principaux, dont la logique peut être expliquée au régulateur.
\item Par les variables qui rendaient chaque alerte inhabituelle, regroupées par type de compte, face aux verdicts des règles, avec les résultats
réinjectés comme étiquettes.
\item La calibration de ses probabilités filtrées (la fréquence de la tension lorsqu'il annonce 0.8), la dérive des paramètres d'une réestimation à l'autre, et
la vraisemblance hors échantillon.
\item L'arbre fixe la hiérarchie de l'allocation~; les groupes de bruit font sauter les pondérations à chaque réestimation et ajoutent de la rotation
sans diversification.
\item Les cas confirmés issus des enquêtes~: lents, peu nombreux, biaisés vers ce que les règles détectent déjà.
\item Ce que disent les données de la veille, et rien de plus.
\end{enumerate}
\end{solution}

\begin{solution}{iq:ml:clustering-regimes-and-anomaly-detection:1}
Filtrage~: la probabilité de l'état sachant les observations jusqu'à $t$, disponible en $t$. Lissage~: sachant toutes les observations, y compris
après $t$~; préférable pour décrire l'historique, mais source de biais d'anticipation (look-ahead bias) s'il sert à décider.
\iqlookfor{les ensembles d'information et le biais d'anticipation.}
\end{solution}

\begin{solution}{iq:ml:clustering-regimes-and-anomaly-detection:2}
Retirer le marché (et éventuellement les facteurs de style), calculer les corrélations sur une fenêtre assez longue pour que les effets apparaissent, partitionner sur la
distance de corrélation, et vérifier la stabilité sur des sous-échantillons et par rapport aux classifications connues.
\iqlookfor{retrait des facteurs, une distance et un contrôle de stabilité.}
\end{solution}

\begin{solution}{iq:ml:clustering-regimes-and-anomaly-detection:3}
EM pour les \omterm{def:ml:clustering-regimes-and-anomaly-detection:hmm}{modèles de Markov cachés}~: l'étape E exécute le forward-backward pour obtenir la probabilité a posteriori de chaque état et de chaque transition~; l'étape M
réestime les probabilités de transition et les paramètres d'émission comme moyennes pondérées par les probabilités a posteriori~; répéter jusqu'à ce que la vraisemblance cesse
de croître.
\iqlookfor{forward-backward, pondérations a posteriori, réestimation, monotonie de la vraisemblance.}
\end{solution}

\begin{solution}{iq:ml:clustering-regimes-and-anomaly-detection:4}
Les arbres aléatoires coupent sur des variables tirées au hasard, à des valeurs aléatoires~; les anomalies demandent moins de coupures pour être isolées. Angles morts~: les anomalies qui
ne sont pas rares (groupes de manipulateurs), les anomalies visibles seulement dans des combinaisons que les coupures aléatoires séparent rarement, et
les variables non pertinentes qui ajoutent du bruit.
\iqlookfor{le mécanisme et au moins deux angles morts.}
\end{solution}

\begin{solution}{iq:ml:clustering-regimes-and-anomaly-detection:5}
Si les probabilités de régime ont été filtrées et estimées uniquement sur des données passées~; si le nombre de régimes et les paramètres ont été
choisis sur tout l'échantillon~; la significativité du résultat~; et sa stabilité sur les sous-périodes.
\iqlookfor{le lissage et l'ajustement sur tout l'échantillon comme premiers suspects.}
\end{solution}

\begin{solution}{iq:ml:clustering-regimes-and-anomaly-detection:6}
Des règles pour les schémas connus, complétées par des scores non supervisés sur des variables construites~; des seuils fixés selon la capacité des analystes~; une évaluation sur
scénarios insérés et cas leurres, sur les rares cas confirmés, et par les décisions des analystes au fil du temps~; une piste d'audit de chaque
alerte et chaque décision.
\iqlookfor{règles et scores combinés, seuils fondés sur la capacité, évaluation sur cas insérés, boucle de retour.}
\end{solution}
